\documentclass[11pt,compress,t,notes=noshow, aspectratio=169, xcolor=table]{beamer}

\usepackage{../../style/lmu-lecture}
% Defines macros and environments
\input{../../style/common.tex}

\title{Applied Machine Learning}
% \author{LMU}
%\institute{\href{https://compstat-lmu.github.io/lecture_iml/}{compstat-lmu.github.io/lecture\_iml}}
\date{}

\begin{document}


\begin{frame}{Calibration versus Discrimination}

\begin{itemize}
\setlength\itemsep{2em}
    \item Calibration indicates the agreement between estimated probabilities and actual observed probabilities (through observed events).
    \item Discrimination is a model's ability to correctly separate observations into specific groups.
    \item Models can over-/underestimate probabilities (poor calibration) and still separate classes (good discrimination) or vice versa.
    \item Consider a model that predicts a patient's disease risk, where a patient is classified as a high-risk patient if disease risk exceeds 20\%. Patient A has a true risk of 5\% while patient B has a true risk of 30\%. The model predicts roughly the same risk with 19\% for A and 21\% for B. It is poorly calibrated but still separates both patients accurately into low-risk and high-risk groups (good discrimination).
    
\end{itemize}

\end{frame}


\endlecture
\end{document}
